\section{实验设计}\label{sec:setup}

\subsection{三档模型}

为了把“结构先验”与“学习”分开，我们设三档，它们共享同一个任务（给定图与源点 $s$，预测到各点的距离）：

\begin{enumerate}[leftmargin=1.8em]
\item \textbf{L0（恒等 min-sum）}：定理~\ref{thm:equiv} 的模型，取 $g(w)=w$、不做任何训练、推理时迭代 128 轮。它\textbf{就是} Bellman-Ford，参数量为 0。
\item \textbf{L1（训练后的 min-sum）}：把边权变换换成可学习的 $g_\theta(w)=\mathrm{softplus}(a)\cdot w+b$（2 个参数，$a=b=0$ 初始化），用监督损失训练 40 轮，推理时同样迭代 128 轮。
\item \textbf{L2（mean 聚合 GNN）}：黑盒基线。每层独立参数的 mean 聚合 + MLP 更新，隐层 64，深度 8，参数量 $104{,}321$。
\end{enumerate}

需要强调：L1 与 L2 的差异是\textbf{结构先验 + 迭代预算 + 参数量}三者的联合差异，不是受控消融。
真正干净的因果对照只有 L0 与 L1——它们同为 128 轮、只差 2 个参数。

\subsection{数据}

训练集：200 张 ER 图（$n\in[12,50]$，平均度 3.0），每张取 2 个源点，共 400 个样本。
测试集六个：同分布 ER、跨生成器的 BA / 网格 / 随机块模型（SBM）、以及\emph{规模外推}（同生成器、同平均度，$n$ 放大到 100 与 200）。
所有图由固定种子生成，同一种子两次生成的边表与权重逐字节相同。

评测目标统一取 $(s,t)$ 这一对：$t$ 取 $s$ 的\emph{跳数最远点}（并列时取加权最远）。
这样选是为了避免 $s$--$t$ 只差一两跳、深度拉不开差距——早期版本按加权最远选 $t$，实测 12.4\% 的实例 $s$--$t$ 只差 1--2 跳，等于把深度这个变量废掉了。

\subsection{三轴指标}

\textbf{准确率}：平均绝对误差（MAE）、命中率（$|\hat d-d|\le0.5$ 的比例）、以及\textbf{最坏误差}。
最坏误差必须单列——第~\ref{sec:results}~节会看到均值把灾难性错误藏起来了。

\textbf{鲁棒性}：对测试图施三种扰动（权重乘 $1\pm20\%$ 的随机因子、随机删 10\% 的边、规模外推），测扰动后的 MAE。
对被删边导致图断开的情形，逐条把边加回直到连通，并记录实际删边比例。
另外单列“任务本身变了多少”（即真实距离的相对变化），用来确认扰动不是摆设。

\textbf{可解释性}：见第~\ref{sec:interp}~节，这里先说明为什么不能直接比 Fidelity。
Fidelity$^+$ 的定义是“只保留解释边之后，预测还保得住多少”。对结构臂它\textbf{恒等于 1}——argmin 链就是最短路，只留这些边当然等价——所以它是恒等式，不是测量。
因此我们另用三个同尺度、有真值的指标：解释边落在真实最短路上的比例（精度）、解释边能否连成 $s$--$t$ 路、以及把解释边权重放大 10 倍后\emph{真实}距离的变化。

\subsection{统计方法}

每个配置跑 8 个种子，报告均值、中位数与逐种子值。
跨臂比较用 Mann-Whitney U（非配对）并给 bootstrap 95\% 置信区间。
\textbf{不用配对检验}：8 个种子是 8 个不同的模型，第 $i$ 号种子在两个臂之间没有自然配对关系，按种子下标配对会人为压低 $p$ 值。

\subsection{负控制与可复现}

只测“能跑通”的断言等于没测。我们为每个验收脚本配了桩实现（例如把 Dijkstra 换成恒返回 0），要求脚本\textbf{必须失败}：
9 个脚本在桩实现下全部正确返回非零。
此外，桩实现会照常写入结果文件，因此负控制运行前会备份 \texttt{results/} 并在结束后恢复——否则会把真实结果覆盖成假数据（我们确实这么被坑过一次）。

复现只需一条命令序列，见仓库根目录的 \texttt{README.md}；全部实验在纯 CPU 上完成，总耗时约 10 分钟。
